% Source: Wald GR, physical PDF pp. 17--20
%         (printed pp. 4--7).
% Coverage: Section 1.2, Space and Time in Prerelativity Physics and in Special
%           Relativity; equations (1.2.1)--(1.2.8).
% Figures/tables/footnotes: Figures 1.1--1.3; no tables or footnotes.
% Status: source-reviewed and compile-clean.
% Uncertainties: none.

\subsection{Space and Time in Prerelativity Physics and in Special Relativity}\label{sec:1.2}

Perhaps the greatest obstacle to understanding the theories of special and general
relativity arises from the difficulty in realizing that a number of previously held basic
assumptions about the nature of space and time are simply wrong. We begin,
therefore, by spelling out some key assumptions about space and time. In both the
past and modern viewpoints, space and time have at least the following structure in
common. We can consider space and time (\(\equiv\) spacetime) to be a continuum
composed of \emph{events}, where each event can be thought of as a point of space at
an instant of time. Furthermore, all events (or, at least, all events in a sufficiently
small neighborhood of a given event) can be uniquely characterized by four numbers:
in ordinary language, three numbers for the spatial position and one for the time. As
will be discussed in chapter 2, a mathematically precise statement of these ideas is
that spacetime is a four-dimensional manifold.

However, prior to relativity theory it was believed that spacetime had the following
additional structure: Given an event \(p\) in spacetime, there is a natural,
observer-independent notion of events occurring ``at the same time'' as \(p\). More
precisely, given two events \(p\) and \(q\), one of the following three mutually
exclusive possibilities must hold: (1) It is possible, in principle, for an observer or
material body to go from event \(q\) to event \(p\), in which case one says \(q\) is to
the past of \(p\). (2) It is possible to go from \(p\) to \(q\), in which case one says
\(q\) is to the future of \(p\). (3) It is impossible, in principle, for an observer or
material body to be present at both events \(p\) and \(q\). In prerelativity physics,
events in the third category are assumed to form a three-dimensional set and define
the notion of simultaneity with \(p\), as is illustrated in Figure~\ref{fig:1.1}.

% Source: Wald GR, physical PDF p. 18 (printed p. 5).
% Coverage: Figure 1.1, causal structure of spacetime in prerelativity physics.
% Figures/tables/footnotes: Figure only; no tables or footnotes.
% Status: source-reviewed and compile-clean.
% Uncertainties: none.

\begingroup
\renewcommand{\thefigure}{1.1}
\begin{figure}[htbp]
  \centering
  \begin{tikzpicture}[>=Stealth, line cap=round, line join=round]
    \draw[->, thick] (0,4.40) -- (0,5.05) node[above] {time};
    \draw[->, thick] (0,4.40) -- (0.85,4.40) node[right] {space};
    \draw[->, thick] (0,4.40) -- (-0.34,4.06);

    \draw[thick] (0.65,2.20) -- (5.65,2.20) -- (6.90,3.75) --
      (1.90,3.75) -- cycle;
    \fill (3.40,2.94) circle (1.7pt) node[left=4pt] {\(p\)};
    \node at (4.15,4.35) {Future};
    \node at (3.10,1.60) {Past};
    \node[anchor=west] at (6.25,4.10) {Simultaneous};
    \draw[->, thick, bend right=18] (6.36,3.97) to (5.88,3.10);
  \end{tikzpicture}
  \caption{A diagram showing the causal structure of spacetime in prerelativity
  physics. Given an event \(p\), all other events in spacetime either are to the future
  of \(p\), to the past of \(p\), or simultaneous with \(p\). The simultaneous events
  form a three-dimensional surface in spacetime.}
  \label{fig:1.1}
\end{figure}
\endgroup


The belief that the causal structure of spacetime has the character shown in
Figure~\ref{fig:1.1} turns out to be wrong. In special relativity theory the above
classification of the causal relationships between events still holds. The crucial
difference is that events in category (3) form more than a three-dimensional set; the
causal relation between \(p\) and other events has the structure sketched in
Figure~\ref{fig:1.2}. The events in category (3) can be further subdivided as follows:
(i) Events that lie on the boundary of the set of points to the future of \(p\). These
events cannot be reached by a material particle starting at \(p\) but can be reached by
a light signal emitted from \(p\). They form the ``future light cone'' of \(p\) (a
three-dimensional set). (ii) Events on the past light cone of \(p\), defined similarly.
(iii) Events in category (3) which are on neither the past nor the future light cone.
These events are said to be spacelike related to \(p\) and comprise a four-dimensional
set.

A key fact closely related to the above is that in special relativity there is no notion
of absolute simultaneity; there are no absolute three-dimensional surfaces in
spacetime as in Figure~\ref{fig:1.1}. As we shall see below, an observer still can
define a notion of which events occur ``at the same time'' as a given event---thus
defining a three-dimensional surface in spacetime---but the notion he gets depends
upon his state of motion. (On the other hand, the light cones of Fig.~\ref{fig:1.2} are
absolute surfaces.) The notion that there is absolute simultaneity is a deeply
ingrained one. The fact that there is no such notion is one of the most difficult ideas
to adjust to in the theory of special relativity.

% Source: Wald GR, physical PDF p. 18 (printed p. 5).
% Coverage: Figure 1.2, causal structure of spacetime in special relativity.
% Figures/tables/footnotes: Figure only; no tables or footnotes.
% Status: source-reviewed and compile-clean.
% Uncertainties: none.

\begingroup
\renewcommand{\thefigure}{1.2}
\begin{figure}[htbp]
  \centering
  \begin{tikzpicture}[>=Stealth, line cap=round, line join=round]
    \draw[thick] (0,0) -- (-1.80,2.20);
    \draw[thick] (0,0) -- (1.80,2.20);
    \draw[thick] (0,0) -- (-1.80,-2.20);
    \draw[thick] (0,0) -- (1.80,-2.20);
    \draw[thick] (0,2.20) ellipse [x radius=1.80, y radius=0.34];
    \draw[thick] (0,-2.20) ellipse [x radius=1.80, y radius=0.34];
    \fill (0,0) circle (1.7pt) node[left=4pt] {\(p\)};

    \node at (0,3.05) {Future};
    \draw[->, thick, bend right=14] (0.60,2.86) to (0.48,2.30);
    \node at (0,-3.05) {Past};
    \draw[->, thick] (0,-2.78) -- (0,-2.43);

    \node[anchor=west] at (2.35,1.05) {Future light cone};
    \draw[->, thick] (2.20,0.94) -- (1.18,1.38);
    \node[anchor=west] at (2.20,-1.35) {Past light cone};
    \draw[->, thick] (2.05,-1.27) -- (1.12,-1.17);
    \node[anchor=east] at (-2.15,0.10) {Spacelike related};
    % Point into the exterior of the cone, not at its vertex p.
    \draw[->, thick] (-2.00,0.02) -- (-1.05,0.36);
    \draw[->, thick] (-2.00,-0.02) -- (-1.08,-0.38);
  \end{tikzpicture}
  \caption{A diagram showing the causal structure of spacetime in special relativity.
  The ``light cone'' of \(p\) rather than a ``surface of simultaneity'' with \(p\) now
  plays a fundamental role in determining the causal relationship of \(p\) to other
  events.}
  \label{fig:1.2}
\end{figure}
\endgroup


In special relativity (as in prerelativity physics) one has the notion of inertial,
``nonaccelerating'' motion, namely the motion a material body would undergo if
subjected to no external forces. An inertial observer can label the events of spacetime
in the following manner. He can build himself a rigid frame and label the grid points
of the frame with the Cartesian coordinates \(x, y, z\) of the (assumed Euclidean)
geometry of the frame. He can then have a clock placed at each grid point and can
synchronize each clock with his by a symmetrical procedure, e.g., by making sure
that a given clock and his give the same reading when they receive a signal sent out
in a symmetrical manner by an observer stationed halfway between the two. (Because
the causal structure of spacetime is that of Figure~\ref{fig:1.2}, not
Figure~\ref{fig:1.1}, synchronization is not a trivial issue.) The observer may carry
the grid, complete with synchronized clocks, in a nonrotating manner. Each event in
spacetime can now be labeled with the coordinates \(x, y, z\) of the grid point at
which the event occurred and the reading \(t\) of the (synchronized) clock at that
event. The labels \(t, x, y, z\) assigned to events in this manner are referred to as
global inertial coordinates.

If two such inertial observers go through this procedure, one may compare the
coordinate labels they assign to events. In prerelativity physics (where the same
labeling procedure works, the only difference being that clock synchronization is
trivial), if observer \(O\) labels an event \(p\) with coordinates \(t, x, y, z\) and
\(O'\) moves with velocity \(v\) in the \(x\)-direction, passing observer \(O\) at the
event labeled by \(t = x = y = z = 0\), the coordinate labels that \(O'\) assigns to
event \(p\) are

\begin{equation}
  t' = t,
  \tag{1.2.1}\label{eq:1.2.1}
\end{equation}
\begin{equation}
  x' = x-vt,
  \tag{1.2.2}\label{eq:1.2.2}
\end{equation}
\begin{equation}
  y' = y,
  \tag{1.2.3}\label{eq:1.2.3}
\end{equation}
\begin{equation}
  z' = z.
  \tag{1.2.4}\label{eq:1.2.4}
\end{equation}

In special relativity, however, the labeling by \(O'\) will be related to that of \(O\)
by a Lorentz transformation,

\begin{equation}
  t' = \bigl(t-vx/c^2\bigr)/\bigl(1-v^2/c^2\bigr)^{1/2},
  \tag{1.2.5}\label{eq:1.2.5}
\end{equation}
\begin{equation}
  x' = \bigl(x-vt\bigr)/\bigl(1-v^2/c^2\bigr)^{1/2},
  \tag{1.2.6}\label{eq:1.2.6}
\end{equation}
\begin{equation}
  y' = y,
  \tag{1.2.7}\label{eq:1.2.7}
\end{equation}
\begin{equation}
  z' = z,
  \tag{1.2.8}\label{eq:1.2.8}
\end{equation}

where \(c\) is the speed of light. Equation~\eqref{eq:1.2.5} shows that the notion
of simultaneity determined by \(O\) (namely, \(t = \text{constant}\)) differs from
that determined by \(O'\) (\(t' = \text{constant}\)), as illustrated in
Figure~\ref{fig:1.3}.

% Source: Wald GR, physical PDF p. 20 (printed p. 7).
% Coverage: Figure 1.3, disagreement over simultaneity between inertial observers.
% Figures/tables/footnotes: Figure only; no tables or footnotes.
% Status: source-reviewed and compile-clean.
% Uncertainties: none.

\begingroup
\renewcommand{\thefigure}{1.3}
\begin{figure}[htbp]
  \centering
  \begin{tikzpicture}[>=Stealth, line cap=round, line join=round]
    \draw[->, thick] (-3.15,2.95) -- (-3.15,3.60) node[above] {\(t\)};
    \draw[->, thick] (-3.15,2.95) -- (-2.35,2.95) node[right] {\(x\)};

    \draw[thick] (0,-2.50) -- (0,2.45);
    \draw[thick] (-2.55,0) -- (3.55,0);
    \draw[dashed, thick] (-2.10,-1.42) -- (2.35,1.50);
    \draw[dashed, thick] (-1.15,-2.70) -- (1.20,2.70);
    \fill (0,0) circle (1.8pt) node[above left=2pt] {\(p\)};

    \node[anchor=south] at (-0.55,1.85) {\(O\)};
    \draw[->, thick, bend left=13] (-0.46,1.75) to (-0.08,1.50);
    \node[anchor=south] at (0.65,1.90) {\(O'\)};
    \draw[->, thick, bend right=15] (0.67,1.80) to (0.54,1.52);

    \node[anchor=west, align=left] at (2.20,0.88) {Surface of\\simultaneity for \(O'\)};
    \draw[->, thick] (2.15,0.75) -- (1.78,1.12);
    \node[anchor=west, align=left] at (1.38,-0.98) {Surface of\\simultaneity for \(O\)};
    \draw[->, thick] (1.32,-0.82) -- (1.22,-0.12);
  \end{tikzpicture}
  \caption{A spacetime diagram illustrating the fact that in special relativity the
  inertial observers \(O\) and \(O'\) disagree over the definition of simultaneity with
  event \(p\).}
  \label{fig:1.3}
\end{figure}
\endgroup

